% Alternative draft aligned with the three RQs in sections/1_intro.tex.
% Use instead of sections/2_related.tex, not alongside it.
% Additional citation entries are in related_work_additions.bib.
% To integrate later, use \bibliography{aaai2026,related_work_additions}.
\section{Related Work}
\label{sec:related-work}

We connect three strands of research: the reception of minority judgments,
the moral reasons exchanged in online discussions, and interventions that
shape readers' exposure to different views. Together, they motivate
examining dissent through community responses, the content of replies,
and readers' consideration of alternative reasons.

\subsubsection*{Minority Reception and Moral Direction}
%\label{subsec:related-reception}
Social influence involves both treating others' judgments as evidence and
seeking social acceptance \cite{deutsch1955normative}. Although individuals
may conform to unanimous majorities \cite{asch1956independence}, consistent
minorities can also influence majority responses
\cite{moscovici1969minority}. Such influence is more evident on private,
indirect measures than in public responses \cite{wood1994minority}.
Reception and influence therefore need to be distinguished: a minority
position may receive little public approval while still affecting how
others approach an issue. Group research likewise shows that dissenters
can attract both communication and rejection
\cite{schachter1951deviation,wesselmann2014revisiting}, with rejection
varying with the relative size of the dissenting faction
\cite{tata1996proportionate}.

Online feedback makes the distinction between evaluation and attention
especially relevant. Prior ratings can shape subsequent evaluations
\cite{muchnik2013social,glenski2017rating}, while Reddit voting also
reflects local norms and practices \cite{graham2021sociomateriality}.
Negatively rated comments can attract more replies and larger reply trees
\cite{davis2021emotional}, and different response behaviors need not
indicate the same reaction \cite{kalch2017replying}. Within AITA,
\citet{mellody2025consensus} finds that the relative valuation of agreement
and dissent varies nonlinearly with the strength of the developing
consensus: dissent is favored at intermediate levels, whereas no
significant preference emerges at very strong consensus. This work makes
agreement context central to understanding reception. Our analysis
compares discussions by their final judgment composition and extends the
comparison from scores to reply counts and depth, while distinguishing
the moral direction of dissent.

Moral direction provides a further distinction among dissenting
judgments. Blame depends on assessments of norm violations, causality,
intentionality, and reasons or preventability \cite{malle2014blame}.
It is also socially regulated: people who assign blame may themselves
face demands for justification or criticism of unwarranted condemnation
\cite{malle2022socially}. Thus, opposing a majority that defends an original
poster calls for condemnation, whereas opposing a blaming majority
resists it. Both positions are numerically minoritarian, but they make
different moral claims. This distinction motivates examining blaming and
defending dissent separately within high-agreement discussions.

Anticipated reception can also shape which views are expressed and how
firmly they are stated. Accountability can encourage accommodation to an
audience, consideration of alternatives, or bolstering of prior
commitments \cite{tetlock1989strategies}. Research on the spiral of silence
links expression to opinion climate and attitude certainty
\cite{matthes2010spiral}, while anticipated personal attacks are associated
with willingness to express minority views \cite{neubaum2018sanctions}.
Expressions of doubt also receive fewer recommendations in online news
discussions \cite{evans2023doubt}. These findings motivate examining
expressed certainty alongside reception in \textbf{RQ1}, while recognizing
that linguistic certainty need not reveal private confidence and that
aggregate patterns may reflect selective participation.

\subsubsection*{Moral Reasons in Online Exchanges}
%\label{subsec:related-reasons}
AITA research characterizes the topics and judgments found in everyday
moral dilemmas \cite{nguyen2022mapping,yudkin2025large}, demographic and
narrative correlates of evaluations
\cite{decandia2022socialnorms,giorgi2023author}, and the situational and
relational contexts of moral disagreement \cite{shi2025moral}. Work on
growing conversation networks \cite{goglia2024structure} and judgment and
language patterns around AITA's judgment disclosure
\cite{goglia2026influence} further situates judgments within unfolding
discussions. Together, these studies characterize moral evaluations in
relation to both the narrated situation and the conversational setting.

Understanding disagreement also requires examining the reasons supporting
judgments. Social Chemistry 101 represents situation-specific norms as
natural-language rules of thumb \cite{forbes2020social}, and Value
Kaleidoscope represents potentially conflicting values, rights, and duties
\cite{sorensen2024value}. More directly, \citet{xi2024morality} extract and
categorize rationale elements in AITA comments and relate their
distributions to social context. These approaches provide ways to
characterize moral content within contributions. A complementary question
is how that content develops between contributions: whether a reply
introduces a moral reason absent from the comment it addresses.

Minority-influence research suggests why these exchanges may matter even
when dissent is rejected. Dissent can encourage divergent thinking
\cite{nemeth1986differential,nemeth1987divergent}, and groups with a history
of minority dissent can exhibit greater cognitive complexity despite
social rejection \cite{curseu2012socially}. Deliberative communication
research similarly distinguishes agreement from the ability to articulate
reasons for different positions. \citet{cappella2002argument} develop the
concept of \emph{argument repertoire}, encompassing reasons supporting
one's own and opposing views; \citet{price2002disagreement} associate
exposure to disagreement with a greater ability to offer reasons for
opposing positions. These findings motivate attention to reasons beyond
whether an exchange produces agreement. In \textbf{RQ2}, we examine this
at the level of direct replies, measuring support and challenge
separately from additional moral reasons. This allows both supportive and
challenging replies to extend their parent's reasoning without treating
addition as evidence of argument quality or subsequent uptake.

\subsubsection*{Information Exposure and Reader Responses}
%\label{subsec:related-exposure}
The presence of different views in a discussion does not ensure that
readers encounter them. Research on Facebook distinguishes the diversity
available through social connections, the content displayed by ranking
algorithms, and the content users choose to consume
\cite{bakshy2015exposure}. Interventions can therefore target different
parts of the exposure process. \citet{munson2010presenting} vary the mix
and presentation of supporting and challenging political opinions,
identifying differences between diversity-seeking and challenge-averse
readers. Such differences suggest that providing diverse views and
securing readers' willingness to engage with them are separate design
questions.

Interventions can also encourage reflection on the sources of exposure.
The Social Mirror \cite{gillani2018mirror} system
lets Twitter users locate themselves in a network
visualization and, in one treatment, receive recommendations for accounts
with opposing political views. Their experiment distinguishes beliefs
about network homogeneity, the diversity of followed accounts, and the
political alignment of shared URLs. Differences between treatments in
network beliefs and subsequent following behavior were not accompanied
by detectable changes in the average alignment of shared URLs. This
distinction highlights the need to evaluate what an intervention changes,
rather than assuming that reflection, exposure, and expression move
together.

At the level of discussion interfaces, Opinion Space supports navigation
among comments from different perspectives and increases engagement and
respect for comments relative to a chronological list, without increasing
the diversity of comments read \cite{faridani2010opinion}. ConsiderIt
organizes public deliberation around shared pro/con reasons, enabling
participants to consider arguments contributed by people with different
positions \cite{kriplean2012considerit}. Comment-ranking diversification
also seeks to include less redundant content among highly visible
comments \cite{northcutt2017diversification}. Complementing these systems,
\citet{dutceac2022cueing} experimentally compare pairs of agreeing and
disagreeing Facebook comments, finding different effects on visual
attention and reported intentions to read or share news. These studies
connect the selection and presentation of comments to distinct forms of
reader engagement.

Field experiments further show that changing exposure does not have a
uniform effect on attitudes. Exposure to opposing political elites on
Twitter increased conservatism among Republican participants in
\citet{bail2018exposure}, whereas counter-attitudinal news subscriptions
on Facebook reduced negative attitudes toward the opposing party in
\citet{levy2021social}. Reducing exposure to like-minded sources changed
Facebook feeds without detectable effects on several preregistered
attitudinal outcomes \cite{nyhan2023likeminded}. More targeted reranking
of content expressing partisan animosity and antidemocratic attitudes
altered affective polarization on X \cite{piccardi2025reranking}.
Collectively, these findings motivate specifying both the content being
surfaced and the outcome being evaluated.

\textbf{RQ3} brings these questions to everyday moral judgment by comparing
two majority threads with one majority and one minority thread. The
comparison makes disagreement available within a single discussion and
examines readers' judgment certainty, reported moral reasons, and
perceived comment influence. Unlike interventions defined by a reader's
prior political position, minority status here is relative to the
discussion's judgment distribution. Our focus is therefore on readers'
responses to the combined presentation of thread content, replies, and
scores, including their reported consideration of reasons beyond those
supporting their final judgment.
